\documentclass[12pt]{article}
\usepackage{amsmath,amssymb,bm}
\usepackage[hidelinks]{hyperref}
\usepackage[nameinlink]{cleveref}

\newcommand{\E}{\mathbb{E}}

\begin{document}

% Dummy anchor so \Cref{eq:main} compiles in isolation
\begin{equation}
\label{eq:main}
Y^* = \alpha + X\eta + \gamma + \sum_{\tau} D_\tau(\Delta\alpha_\tau + X\Delta\eta_\tau + \Delta\gamma_\tau).
\end{equation}

\subsubsection{Why pooling across school types biases logit estimates}
\label{sec:pooling_schooltype_bias}
In several applications it is tempting to pool all schools and impose common slope parameters (including the pandemic-year interaction slopes) across heterogeneous \emph{school types}. This restriction is innocuous in linear models under standard exogeneity assumptions, but in nonlinear likelihood models it generally produces \emph{misspecification} and hence inconsistent slope estimates when the conditional response function differs systematically by type.

Let $s\in\mathcal{S}$ denote school type (e.g.\ comprehensive/grammar/independent) and let $s(j)$ denote the type of school $j$. Consider the type-heterogeneous logit analogue of \Cref{eq:main} in which the index coefficients are allowed to differ by type:
\begin{equation}
\label{eq:main_type}
\Pr\!\left(Y_{i,j,k,t}=1 \mid X_i, j, k, t\right)
=
F\!\Big(z_{i,j,k,t}\big(\theta_{s(j)}\big)\Big),
\end{equation}
where $\theta_{s}$ collects all parameters relevant for that type (including the baseline mapping from $X_i$ to grades and the pandemic-year interaction terms). To make the argument transparent, write a generic regressor vector as $x_{i,j,k,t}$ and a generic slope subvector as $\beta_s$ so that, suppressing fixed effects for notational clarity,
\(
\Pr(Y=1\mid x,s)=F(x'\beta_s).
\)

\vspace{0.5em}
\noindent\textbf{Proposition 1 (Pseudo-true pooled logit and slope bias).}
Suppose the data are generated by $\Pr(Y=1\mid x,s)=F(x'\beta_s)$ with at least two school types $s$ such that $\beta_s\neq \beta_{s'}$ on a set of positive probability, and let $\widehat\beta^{\,P}$ denote the MLE from the \emph{pooled} logit that (incorrectly) imposes a common slope $\beta$ (i.e.\ estimates $F(x'\beta)$). Under standard regularity conditions for misspecified MLE, $\widehat\beta^{\,P}\overset{p}{\to}\beta^{P}$, where the probability limit $\beta^{P}$ solves the pooled population score equation
\begin{equation}
\label{eq:pooled_score}
0
=
\E\!\left[
x\left(
\E[Y\mid x,s]-F(x'\beta^{P})
\right)
\right]
=
\E\!\left[
x\left(
F(x'\beta_{s})-F(x'\beta^{P})
\right)
\right].
\end{equation}
Define the type-weighted mean slope $\bar\beta:=\E[\beta_s]$ and let $f(u):=F(u)\big(1-F(u)\big)$ denote the logit derivative. A first-order expansion of \Cref{eq:pooled_score} around $\bar\beta$ yields
\begin{equation}
\label{eq:beta_bias}
\beta^{P}
=
\bar\beta
+
\beta_{\mathrm{bias}},
\qquad\text{with}\qquad
\beta_{\mathrm{bias}}
\approx
\Big(\E\!\left[f(x'\bar\beta)\,x x'\right]\Big)^{-1}
\E\!\left[f(x'\bar\beta)\,x x'\,(\beta_{s}-\bar\beta)\right].
\end{equation}
Unless $\beta_s$ is constant across types (or heterogeneity is orthogonal to $x$ in the specific sense that the right-hand expectation vanishes), $\beta_{\mathrm{bias}}\neq 0$ and the pooled slope does not recover any type-specific causal response.

\vspace{0.5em}
\noindent\textit{Derivation of \Cref{eq:beta_bias}.}
Using $\E[Y\mid x,s]=F(x'\beta_s)$, the pooled score takes the form in \Cref{eq:pooled_score}. Apply a Taylor expansion around $\bar\beta$:
\(
F(x'\beta_s)\approx F(x'\bar\beta)+f(x'\bar\beta)\,x'(\beta_s-\bar\beta)
\)
and
\(
F(x'\beta^{P})\approx F(x'\bar\beta)+f(x'\bar\beta)\,x'(\beta^{P}-\bar\beta).
\)
Plugging into \Cref{eq:pooled_score} and collecting terms implies
\(
\E\!\left[f(x'\bar\beta)\,x x'\right](\beta^{P}-\bar\beta)
\approx
\E\!\left[f(x'\bar\beta)\,x x'\,(\beta_s-\bar\beta)\right],
\)
which rearranges to \Cref{eq:beta_bias}.

\vspace{0.5em}
\noindent\textbf{Corollary 1 (Correction required for unit-level treatment effects).}
Let $\theta^{P}$ denote the pseudo-true pooled parameter vector (the probability limit of the pooled MLE under misspecification). Let $\tau_{i,j,k,\tau}(\theta)$ denote the unit-level grading-method effect implied by \Cref{eq:main} for pandemic year $\tau$, i.e.\ the change in success probability induced by switching $D_\tau$ from $0$ to $1$ holding $(X_i,j,k)$ fixed:
\begin{equation}
\label{eq:unit_te_def}
\tau_{i,j,k,\tau}(\theta)
:=
F\!\left(z_{i,j,k,t}\big|_{D_\tau=1}\right)
-
F\!\left(z_{i,j,k,t}\big|_{D_\tau=0}\right).
\end{equation}
If the true data-generating process features type-heterogeneous parameters $\theta_{s(j)}$ as in \Cref{eq:main_type}, then in general
\begin{equation}
\label{eq:unit_te_bias}
\tau_{i,j,k,\tau}\!\left(\theta^{P}\right)
\neq
\tau_{i,j,k,\tau}\!\left(\theta_{s(j)}\right),
\end{equation}
because the unit effect is a nonlinear functional of the full index parameters. Therefore, some \emph{correction} is required to consistently estimate unit-level treatment effects: a sufficient correction is to either (i) estimate \Cref{eq:main} separately within each school type $s$ (so that $\widehat\theta_s\overset{p}{\to}\theta_s$), or (ii) equivalently, estimate a single pooled likelihood that fully saturates type heterogeneity by interacting school-type indicators with all slope terms (including the pandemic-year interactions). Under either correction, $\widehat\tau_{i,j,k,\tau}=\tau_{i,j,k,\tau}(\widehat\theta_{s(j)})$ is consistent for $\tau_{i,j,k,\tau}(\theta_{s(j)})$.

\end{document}
